% !PRINTED 508
\weberpage{508}
%<<CONT
when $a\partial = n$ and $c$ is divisible by 8, an integral algebraic
function of the field $\mathfrak{R}(v,u)$, where
\weq{2}{v = j\left(\frac{c+\partial\omega}{2}\right), \qquad u = j(\omega).}
% [sic] German source prints denominator $2$, conflicting with denominator $a$ in (5) and with the ensuing system (7)
Now let a root of the primitive quadratic equation
\weq{3}{A\omega^2 + B\omega + C = 0,}
with negative discriminant
\weq{4}{D = B^2 - 4AC.}
% [sic] German source prints a full stop mid-sentence; the governing verb follows the display
be substituted in it for $\omega$, and we assume once and for all that $A$ is
positive and relatively prime to $2Dn$, which is no restriction. If now
$v = u$ is to hold, the necessary and sufficient condition for it
is:
\weq{5}{\frac{c+\partial\omega}{a} = \frac{\gamma+\delta\omega}{\alpha+\beta\omega},}
\weq{6}{\alpha\delta - \beta\gamma = 1.}
If one compares (5) with (3), it follows (§ 114) that there must be two
positive integers $x, y$, the first of which may be assumed positive, and which satisfy the conditions:
\weq{7}{\begin{aligned}
\partial\beta &= Ax, & c\beta+\partial\alpha-a\delta &= Bx,\\
c\alpha-a\gamma &= Cx, & -c\beta+\partial\alpha+a\delta &= y,\\
2\partial\alpha &= Bx+y, & & \\
2(c\beta-a\delta) &= Bx-y, & &
\end{aligned}}
and from this, on account of (6):
\weq{8}{4n = y^2 - Dx^2.}
A prime divisor of $\partial$ would, since $A$ is assumed relatively prime to $n$, have to divide $x$ and $y$. It would then also divide
\[
c\alpha - a\gamma \quad\text{and}\quad c\beta - a\delta,
\]
and consequently divide $a, \partial, c$. But since these three numbers have no common divisor, $\partial = 1$ must hold, and from (7) there results:
\weq{9}{\begin{aligned}
\alpha &= \frac{Bx+y}{2}, & \beta &= Ax,\\
c\alpha - n\gamma &= Cx, & c\beta - n\delta &= \frac{Bx-y}{2};
\end{aligned}}
from this it follows that $x$ and $y$ have no odd common divisor, and that, if $x$ and $y$ are even, $\frac{1}{2}(Bx+y)$ must be odd. Since we may moreover assume $x$ positive,
%CONT>>
